\chapter{Einführung in maschinelles Lernen}

\section{Ziel dieses Kapitels}

Dieses Kapitel führt die Grundbegriffe des maschinellen Lernens (machine learning) ein.
Es geht noch nicht darum, einen bestimmten Algorithmus vollständig zu beherrschen.
Stattdessen soll klar werden, was maschinelles Lernen überhaupt leisten soll, welche Arten von Lernproblemen es gibt und warum Wahrscheinlichkeitstheorie (probability theory) im maschinellen Lernen eine zentrale Rolle spielt.

\begin{conceptbox}
Maschinelles Lernen (machine learning) beschäftigt sich mit Systemen, die aus Daten Muster lernen und dieses gelernte Wissen auf neue, noch nicht gesehene Daten anwenden sollen.

Das zentrale Ziel ist also nicht bloß, vorhandene Daten auswendig zu lernen, sondern Generalisierung (generalization).
\end{conceptbox}

Nach diesem Kapitel solltest du folgende Fragen beantworten können:

\begin{itemize}
    \item Was ist der Unterschied zwischen künstlicher Intelligenz (artificial intelligence) und maschinellem Lernen (machine learning)?
    \item Was ist ein Datensatz (data set) und was ist ein Modell (model)?
    \item Was bedeutet Lernen in einem mathematischen Sinn?
    \item Was ist der Unterschied zwischen überwachtem Lernen (supervised learning), halbüberwachtem Lernen (semi-supervised learning), unüberwachtem Lernen (unsupervised learning) und bestärkendem Lernen (reinforcement learning)?
    \item Was ist der Unterschied zwischen Klassifikation (classification) und Regression (regression)?
    \item Warum liefern viele moderne Modelle keine festen Antworten, sondern Wahrscheinlichkeiten?
    \item Was bedeutet Training (training), Testen (testing) und Generalisierung (generalization)?
\end{itemize}

\section{Was soll ein KI-System leisten?}

Ein System der künstlichen Intelligenz (artificial intelligence) soll Aufgaben lösen, die typischerweise als intelligent wahrgenommen werden.
Beispiele sind:

\begin{itemize}
    \item ein Bild einer Ziffer erkennen,
    \item eine E-Mail als Spam oder Nicht-Spam klassifizieren,
    \item den zukünftigen Verkaufspreis einer Wohnung vorhersagen,
    \item ähnliche Patientinnen und Patienten gruppieren,
    \item einem Roboter beibringen, sich in einer Umgebung sinnvoll zu bewegen.
\end{itemize}

Maschinelles Lernen (machine learning) ist ein Teilgebiet davon.
Die Grundidee ist:
Wir schreiben nicht für jede mögliche Situation eine feste Regel.
Stattdessen geben wir dem System Daten und lassen es aus diesen Daten eine Regel, Funktion oder Verteilung lernen.

\begin{definitionbox}
Ein Modell (model) ist im maschinellen Lernen eine mathematische Vorschrift, die Eingaben in Ausgaben umwandelt oder die Struktur der Daten beschreibt.

Typisch schreibt man:
\[
    y = f(x).
\]

Dabei ist \(x\) die Eingabe (input), \(f\) das gelernte Modell und \(y\) die Ausgabe (output).
\end{definitionbox}

Die Eingabe \(x\) kann sehr unterschiedlich aussehen:

\begin{itemize}
    \item ein einzelner Messwert, zum Beispiel Temperatur,
    \item ein Vektor von Messwerten, zum Beispiel Blutdruck, Alter und Puls,
    \item ein Bild, das als sehr großer Vektor von Pixelwerten dargestellt wird,
    \item ein Text, der zuerst in eine numerische Darstellung umgewandelt wird.
\end{itemize}

In den meisten ML-Modellen werden Eingaben als Vektoren geschrieben:
\[
    \vec{x}
    =
    \begin{pmatrix}
        x_1 \\
        x_2 \\
        \vdots \\
        x_D
    \end{pmatrix}
    \in \mathbb{R}^D.
\]

Hier ist \(D\) die Dimension (dimension) der Eingabe.
Bei einem Bild mit vielen Pixeln kann \(D\) sehr groß sein.

\section{Daten, Merkmale und Zielwerte}

Ein Datensatz (data set) besteht aus Beobachtungen.
Eine einzelne Beobachtung wird oft als Datenpunkt (data point) bezeichnet.

\begin{definitionbox}
Ein Datenpunkt (data point) ist eine einzelne Beobachtung.
Er besteht meistens aus einem Eingabevektor \(\vec{x}\) und, falls vorhanden, einem Zielwert \(t\).

Ein überwachter Trainingsdatensatz hat typischerweise die Form
\[
    \mathcal{D}
    =
    \left\{
        (\vec{x}_1,t_1),
        (\vec{x}_2,t_2),
        \dots,
        (\vec{x}_N,t_N)
    \right\}.
\]
Dabei ist \(N\) die Anzahl der Datenpunkte.
\end{definitionbox}

Die Komponenten eines Eingabevektors nennt man Merkmale (features).
Ein Beispiel:

\[
    \vec{x}
    =
    \begin{pmatrix}
        \text{Größe} \\
        \text{Gewicht} \\
        \text{Alter}
    \end{pmatrix}.
\]

In einem echten ML-System würden diese Einträge natürlich als Zahlen dargestellt werden, also zum Beispiel:

\[
    \vec{x}
    =
    \begin{pmatrix}
        178 \\
        72 \\
        23
    \end{pmatrix}.
\]

\begin{notebox}
In mathematischen Herleitungen wird oft einfach \(x\) geschrieben, obwohl eigentlich ein Vektor \(\vec{x}\) gemeint ist.
Das ist eine übliche Abkürzung.
In diesem Skript wird möglichst klar unterschieden, aber in späteren Kapiteln kann \(x\) je nach Kontext auch für einen Vektor stehen.
\end{notebox}

\section{Lernen als Funktionsanpassung}

Eine einfache Sichtweise auf maschinelles Lernen ist:

\begin{center}
\textbf{Wir suchen eine Funktion, die zu den Daten passt und für neue Daten gute Vorhersagen macht.}
\end{center}

Angenommen, wir haben Eingaben \(\vec{x}_n\) und Zielwerte \(t_n\).
Dann suchen wir ein Modell \(f\), sodass

\[
    f(\vec{x}_n) \approx t_n.
\]

Das Zeichen \(\approx\) ist wichtig:
In echten Daten gibt es fast immer Rauschen (noise), Messfehler oder unvollständige Information.
Deshalb erwarten wir meistens nicht, dass das Modell alle Daten perfekt trifft.

\begin{examplebox}
Angenommen, wir wollen die Miete einer Wohnung vorhersagen.
Die Eingabe könnte sein:
\[
    \vec{x}
    =
    \begin{pmatrix}
        \text{Wohnfläche} \\
        \text{Anzahl Zimmer} \\
        \text{Entfernung zum Zentrum}
    \end{pmatrix}.
\]
Der Zielwert \(t\) wäre dann die tatsächliche Miete.

Ein Modell \(f\) soll aus diesen Eingaben eine vorhergesagte Miete berechnen:
\[
    \hat{t} = f(\vec{x}).
\]
Der Wert \(\hat{t}\) ist die Vorhersage (prediction).
\end{examplebox}

\section{Parameter eines Modells}

Viele Modelle enthalten Parameter (parameters).
Diese Parameter bestimmen, wie genau das Modell rechnet.

Ein sehr einfaches Modell für eine eindimensionale Eingabe ist eine Gerade:

\[
    y(x,w_0,w_1) = w_0 + w_1x.
\]

Dabei sind \(w_0\) und \(w_1\) Parameter.
Man kann sie zusammen als Vektor schreiben:

\[
    \vec{w}
    =
    \begin{pmatrix}
        w_0 \\
        w_1
    \end{pmatrix}.
\]

Dann schreibt man das Modell oft als

\[
    y(x,\vec{w}).
\]

\begin{definitionbox}
Lernen (learning) bedeutet in vielen ML-Modellen, geeignete Parameter \(\vec{w}\) aus den Daten zu bestimmen.

Informell:
\[
    \text{Daten}
    \longrightarrow
    \text{Lernalgorithmus}
    \longrightarrow
    \text{gelernte Parameter}.
\]
\end{definitionbox}

In späteren Kapiteln wird diese Idee präziser.
Bei der linearen Regression (linear regression) werden die Parameter zum Beispiel so gewählt, dass ein Fehlermaß (error function) möglichst klein wird.
Bei probabilistischen Modellen werden die Parameter oft über Maximum-Likelihood-Schätzung (maximum likelihood estimation) oder Maximum-a-posteriori-Schätzung (maximum a-posteriori estimation) bestimmt.

\section{Trainingsdaten, Testdaten und Generalisierung}

Ein Modell soll nicht nur die Daten erklären, die es schon gesehen hat.
Es soll auch bei neuen Daten funktionieren.
Deshalb unterscheidet man zwischen Trainingsdaten (training data) und Testdaten (test data).

\begin{definitionbox}
Trainingsdaten (training data) sind Daten, mit denen das Modell gelernt wird.

Testdaten (test data) sind Daten, die beim Lernen nicht verwendet werden und danach zur Bewertung des Modells dienen.

Generalisierung (generalization) bedeutet, dass ein Modell auch auf neuen, bisher unbekannten Daten gute Vorhersagen macht.
\end{definitionbox}

Die zentrale Herausforderung ist:

\[
    \text{gute Leistung auf Trainingsdaten}
    \;\not\Rightarrow\;
    \text{gute Leistung auf neuen Daten}.
\]

Ein Modell kann die Trainingsdaten sehr gut treffen und trotzdem schlecht generalisieren.
Dieses Problem nennt man Überanpassung (overfitting).

\begin{conceptbox}
Ein Modell soll die Struktur in den Daten lernen, nicht bloß die Trainingsdaten auswendig.
Wenn ein Modell auch zufälliges Rauschen (noise) lernt, spricht man von Überanpassung (overfitting).
\end{conceptbox}

\section{Klassifikation und Regression}

Zwei grundlegende Problemtypen im überwachten Lernen sind Klassifikation (classification) und Regression (regression).

\subsection{Klassifikation}

Bei einer Klassifikation (classification) soll das Modell eine Eingabe einer von mehreren Klassen (classes) zuordnen.

\begin{definitionbox}
Bei einer Klassifikation (classification) ist der Zielwert diskret.
Das heißt, er stammt aus einer endlichen oder abzählbaren Menge von Klassen:

\[
    t \in \{1,2,\dots,C\}.
\]

Dabei ist \(C\) die Anzahl der Klassen.
\end{definitionbox}

Beispiele:

\begin{itemize}
    \item Bild einer Ziffer \(\rightarrow\) Klasse \(0,1,\dots,9\),
    \item E-Mail \(\rightarrow\) Spam oder Nicht-Spam,
    \item medizinischer Befund \(\rightarrow\) krank oder gesund,
    \item Foto \(\rightarrow\) Katze, Hund, Auto, Person.
\end{itemize}

Für eine Ziffernerkennung könnte das Modell eine Eingabe \(\vec{x}\) bekommen, die ein Bild beschreibt.
Die Ausgabe ist dann zum Beispiel:

\[
    y(\vec{x}) = 7.
\]

Das bedeutet:
Das Modell entscheidet, dass das Bild zur Klasse \(7\) gehört.

\begin{notebox}
In modernen probabilistischen Klassifikationsmodellen gibt das Modell oft nicht direkt eine Klasse aus, sondern Wahrscheinlichkeiten für alle Klassen:

\[
    p(c=0 \mid \vec{x}) = 0{,}01,
    \quad
    p(c=1 \mid \vec{x}) = 0{,}03,
    \quad
    p(c=2 \mid \vec{x}) = 0{,}91,
    \quad
    \dots
\]

Die endgültige Klasse kann dann als die wahrscheinlichste Klasse gewählt werden.
\end{notebox}

\subsection{Regression}

Bei einer Regression (regression) soll das Modell einen kontinuierlichen Wert vorhersagen.

\begin{definitionbox}
Bei einer Regression (regression) ist der Zielwert kontinuierlich, also typischerweise eine reelle Zahl:

\[
    t \in \mathbb{R}.
\]
\end{definitionbox}

Beispiele:

\begin{itemize}
    \item Wohnungsdaten \(\rightarrow\) Mietpreis,
    \item Wetterdaten \(\rightarrow\) Temperatur morgen,
    \item medizinische Daten \(\rightarrow\) erwarteter Laborwert,
    \item technische Prozessdaten \(\rightarrow\) Produktionsausbeute.
\end{itemize}

Ein Regressionsmodell gibt also nicht eine Klasse wie Katze oder Hund aus, sondern zum Beispiel:

\[
    \hat{t} = 23{,}7.
\]

\begin{examtipbox}
Prüfungsfragen unterscheiden oft zwischen Klassifikation (classification) und Regression (regression).

\begin{itemize}
    \item Diskrete Ausgabe? \(\Rightarrow\) Klassifikation (classification).
    \item Kontinuierliche Ausgabe? \(\Rightarrow\) Regression (regression).
\end{itemize}

Wichtig ist die Art des Zielwerts, nicht die Art der Eingabe.
Ein Bild kann zum Beispiel für Klassifikation verwendet werden, aber auch für Regression, wenn man daraus einen kontinuierlichen Wert vorhersagen will.
\end{examtipbox}

\section{Arten des Lernens}

Maschinelles Lernen kann danach eingeteilt werden, welche Information beim Lernen verfügbar ist.

\subsection{Überwachtes Lernen}

\begin{definitionbox}
Beim überwachten Lernen (supervised learning) bestehen die Trainingsdaten aus Eingaben und bekannten Zielwerten:

\[
    \mathcal{D}
    =
    \{(\vec{x}_1,t_1),\dots,(\vec{x}_N,t_N)\}.
\]

Das Modell lernt, Eingaben \(\vec{x}\) auf Zielwerte \(t\) abzubilden.
\end{definitionbox}

Man kann sich das wie Lernen mit einer Lehrperson vorstellen:
Für jeden Trainingsfall ist bekannt, was die richtige Antwort gewesen wäre.

Typische Aufgaben:

\begin{itemize}
    \item Klassifikation (classification),
    \item Regression (regression).
\end{itemize}

\begin{examplebox}
Bei der Ziffernerkennung besteht ein Trainingsbeispiel aus einem Bild und dem zugehörigen Label:

\[
    (\text{Bild einer 7},\,7).
\]

Das Modell soll aus vielen solchen Beispielen lernen, neue Bilder korrekt zu klassifizieren.
\end{examplebox}

\subsection{Halbüberwachtes Lernen}

\begin{definitionbox}
Beim halbüberwachten Lernen (semi-supervised learning) gibt es einige Datenpunkte mit Labels und viele Datenpunkte ohne Labels.
\end{definitionbox}

Die Daten haben also ungefähr die Form:

\[
    \mathcal{D}
    =
    \underbrace{
        \{(\vec{x}_1,t_1),\dots,(\vec{x}_M,t_M)\}
    }_{\text{gelabelte Daten}}
    \cup
    \underbrace{
        \{\vec{x}_{M+1},\dots,\vec{x}_N\}
    }_{\text{ungelabelte Daten}}.
\]

Das ist praktisch wichtig, weil Labels oft teuer sind.
Zum Beispiel kann man viele medizinische Bilder speichern, aber ihre korrekte Diagnose muss von Expertinnen oder Experten erstellt werden.

\subsection{Unüberwachtes Lernen}

\begin{definitionbox}
Beim unüberwachten Lernen (unsupervised learning) gibt es keine Zielwerte.
Die Daten bestehen nur aus Eingaben:

\[
    \mathcal{D}
    =
    \{\vec{x}_1,\dots,\vec{x}_N\}.
\]

Das Ziel ist, Struktur in den Daten zu finden.
\end{definitionbox}

Typische Aufgaben sind:

\begin{itemize}
    \item Cluster finden (clustering),
    \item Wahrscheinlichkeitsdichten schätzen (density estimation),
    \item Daten in niedrigere Dimensionen projizieren (dimensionality reduction),
    \item Muster oder latente Faktoren finden (latent factors).
\end{itemize}

\begin{examplebox}
Angenommen, wir haben Messwerte von vielen Patientinnen und Patienten, aber keine festen Diagnosen.
Ein Clustering-Algorithmus könnte versuchen, Gruppen ähnlicher Patientinnen und Patienten zu finden.
Diese Gruppen müssen nicht vorher bekannt sein.
\end{examplebox}

\subsection{Bestärkendes Lernen}

\begin{definitionbox}
Beim bestärkenden Lernen (reinforcement learning) lernt ein Agent (agent) durch Interaktion mit einer Umgebung (environment).
Der Agent führt Aktionen (actions) aus und erhält Belohnungen (rewards).
Ziel ist es, langfristig eine möglichst hohe Belohnung zu erreichen.
\end{definitionbox}

Die Grundelemente sind:

\begin{itemize}
    \item Zustand (state),
    \item Aktion (action),
    \item Belohnung (reward),
    \item Strategie (policy),
    \item Umgebung (environment).
\end{itemize}

Anders als beim überwachten Lernen bekommt der Agent nicht für jede Situation die perfekte Antwort gezeigt.
Er muss durch Versuch und Irrtum (trial and error) lernen.

\begin{examplebox}
Ein Spielagent sieht eine Spielsituation, wählt einen Zug und bekommt später eine Belohnung, zum Beispiel \(+1\) für einen Sieg und \(-1\) für eine Niederlage.

Das Schwierige ist:
Ein Zug kann erst viele Schritte später wichtig werden.
Man muss also lernen, welche früheren Aktionen zu späterem Erfolg beigetragen haben.
\end{examplebox}

\begin{notebox}
Bestärkendes Lernen (reinforcement learning) ist wichtig, aber in diesem Skript liegt der Schwerpunkt zunächst auf probabilistischen Grundlagen, Regression, Klassifikation, neuronalen Netzen, Dimensionsreduktion, Clustering, Gaussian Mixture Models und dem EM-Algorithmus.
\end{notebox}

\section{Probabilistische Antworten als optimale Inferenz}

Viele ML-Probleme haben keine absolut sichere Antwort.
Ein Bild kann undeutlich sein, Messwerte können rauschen, Daten können unvollständig sein.
Deshalb ist es oft sinnvoller, nicht nur eine feste Antwort zu geben, sondern Unsicherheit mitzuschätzen.

\begin{conceptbox}
Ein probabilistisches Modell (probabilistic model) gibt nicht nur eine Vorhersage aus, sondern beschreibt Unsicherheit mit Wahrscheinlichkeiten.
\end{conceptbox}

Bei einer Klassifikation könnte ein Modell zum Beispiel ausgeben:

\[
    p(c=0 \mid \vec{x}) = 0{,}01,
    \quad
    p(c=1 \mid \vec{x}) = 0{,}03,
    \quad
    p(c=2 \mid \vec{x}) = 0{,}91,
    \quad
    \dots
\]

Dann wäre Klasse \(2\) die wahrscheinlichste Klasse.
Aber das Modell sagt zusätzlich:
Es ist sich zu etwa \(91\%\) sicher.

\subsection{Warum ist Unsicherheit nützlich?}

Unsicherheit ist nützlich, weil sie bessere Entscheidungen erlaubt.

\begin{itemize}
    \item Bei hoher Sicherheit kann ein System automatisch entscheiden.
    \item Bei niedriger Sicherheit kann ein Mensch nachprüfen.
    \item Bei medizinischen oder sicherheitskritischen Aufgaben ist Unsicherheit besonders wichtig.
    \item Wahrscheinlichkeiten erlauben es, mehrere Informationsquellen mathematisch sauber zu kombinieren.
\end{itemize}

\begin{examplebox}
Ein medizinisches Modell klassifiziert ein Bild als unauffällig oder auffällig.

Fall A:
\[
    p(\text{auffällig} \mid \vec{x}) = 0{,}99.
\]

Fall B:
\[
    p(\text{auffällig} \mid \vec{x}) = 0{,}51.
\]

In beiden Fällen wäre die wahrscheinlichste Klasse auffällig.
Trotzdem sind die Situationen sehr verschieden:
In Fall A ist das Modell sehr sicher, in Fall B fast unsicher.
\end{examplebox}

\section{Ein erstes Klassifikationsbeispiel}

Wir betrachten ein einfaches Klassifikationsproblem.
Ein System bekommt ein Bild einer handgeschriebenen Ziffer.
Die möglichen Klassen sind:

\[
    c \in \{0,1,2,3,4,5,6,7,8,9\}.
\]

Das Bild wird als Eingabevektor \(\vec{x}\) dargestellt.
Wenn das Bild zum Beispiel \(28 \times 28\) Pixel hat, dann enthält es

\[
    28 \cdot 28 = 784
\]

Pixelwerte.
Das Bild kann also als Vektor in \(\mathbb{R}^{784}\) dargestellt werden:

\[
    \vec{x} \in \mathbb{R}^{784}.
\]

Das Modell soll aus \(\vec{x}\) die richtige Ziffer bestimmen.

\begin{definitionbox}
Bei der Ziffernklassifikation ist

\[
    \vec{x} = \text{Bild als Vektor},
    \qquad
    c = \text{Ziffernklasse}.
\]

Ein probabilistischer Klassifikator modelliert dann Wahrscheinlichkeiten der Form

\[
    p(c \mid \vec{x}).
\]
\end{definitionbox}

Die endgültige Entscheidung könnte durch die wahrscheinlichste Klasse getroffen werden:

\[
    \hat{c}
    =
    \arg\max_{c \in \{0,\dots,9\}}
    p(c \mid \vec{x}).
\]

Das bedeutet:
Wir wählen die Klasse \(c\), für die die bedingte Wahrscheinlichkeit \(p(c \mid \vec{x})\) maximal ist.

\begin{notebox}
Das Symbol \(\arg\max\) liefert nicht den maximalen Wert selbst, sondern die Stelle, an der das Maximum erreicht wird.

Wenn \(p(c=7 \mid \vec{x})\) am größten ist, dann ist

\[
    \arg\max_c p(c \mid \vec{x}) = 7.
\]
\end{notebox}

\section{Merkmalsextraktion und Vorverarbeitung}

Rohdaten sind oft nicht direkt ideal für ein ML-Modell.
Deshalb werden Daten häufig vorverarbeitet (preprocessed).

\begin{definitionbox}
Vorverarbeitung (preprocessing) bezeichnet Transformationen der Rohdaten, bevor sie in das Modell eingegeben werden.

Merkmalsextraktion (feature extraction) bedeutet, aus Rohdaten geeignete Merkmale (features) zu berechnen.
\end{definitionbox}

Beispiele:

\begin{itemize}
    \item Bilder werden skaliert oder zentriert.
    \item Messwerte werden normalisiert.
    \item Texte werden in numerische Vektoren umgewandelt.
    \item Sehr hochdimensionale Daten werden auf wenige Dimensionen reduziert.
\end{itemize}

\begin{examtipbox}
Wichtig:
Die gleiche Vorverarbeitung muss auf Trainingsdaten und Testdaten angewendet werden.
Wenn Trainingsdaten anders verarbeitet werden als Testdaten, misst man nicht mehr die echte Modellleistung.
\end{examtipbox}

\section{Modell, Verlustfunktion und Optimierung}

Viele ML-Verfahren lassen sich in drei Bausteine zerlegen:

\begin{enumerate}
    \item Modell (model),
    \item Verlustfunktion (loss function),
    \item Optimierungsverfahren (optimization method).
\end{enumerate}

\subsection{Modell}

Das Modell beschreibt, welche Funktionen überhaupt möglich sind.
Zum Beispiel:

\[
    y(x,\vec{w}) = w_0 + w_1x
\]

für eine Gerade oder ein neuronales Netz (neural network) für eine komplexere Abbildung.

\subsection{Verlustfunktion}

Die Verlustfunktion (loss function) misst, wie schlecht das Modell auf den Trainingsdaten ist.
Eine typische Idee ist:

\[
    \text{großer Fehler}
    \quad\Rightarrow\quad
    \text{großer Verlust}.
\]

Für Regression wird oft ein quadratischer Fehler verwendet:

\[
    E(\vec{w})
    =
    \frac{1}{2}
    \sum_{n=1}^{N}
    \left(
        y(\vec{x}_n,\vec{w}) - t_n
    \right)^2.
\]

Der Faktor \(\frac{1}{2}\) ist mathematisch praktisch, weil er sich beim Ableiten mit dem Faktor \(2\) aus dem Quadrat kürzt.

\subsection{Optimierung}

Beim Training sucht man Parameter, die den Verlust möglichst klein machen:

\[
    \vec{w}^{\,*}
    =
    \arg\min_{\vec{w}} E(\vec{w}).
\]

\begin{definitionbox}
Training (training) bedeutet oft:

\[
    \text{Finde Parameter, die eine Verlustfunktion minimieren.}
\]

Mathematisch schreibt man:

\[
    \vec{w}^{\,*}
    =
    \arg\min_{\vec{w}} E(\vec{w}).
\]
\end{definitionbox}

Später werden wir sehen, dass dieses Optimierungsproblem bei manchen Modellen direkt gelöst werden kann, während man bei anderen Modellen iterative Verfahren wie Gradientenabstieg (gradient descent) verwendet.

\section{Überanpassung und Unteranpassung}

Ein Modell kann auf verschiedene Arten schlecht sein.

\subsection{Unteranpassung}

Unteranpassung (underfitting) bedeutet, dass das Modell zu einfach ist, um die Struktur der Daten zu erfassen.

\begin{examplebox}
Wenn die Daten ungefähr einer gekrümmten Kurve folgen, aber man nur eine Gerade anpasst, kann die Gerade die Struktur nicht gut treffen.
Das Modell ist dann zu unflexibel.
\end{examplebox}

\subsection{Überanpassung}

Überanpassung (overfitting) bedeutet, dass das Modell zu stark an die Trainingsdaten angepasst ist und dabei auch Rauschen (noise) lernt.

\begin{examplebox}
Ein sehr komplexes Polynom kann durch jeden Trainingspunkt laufen.
Das sieht auf den Trainingsdaten perfekt aus, kann aber zwischen den Datenpunkten stark oszillieren und auf neuen Daten schlecht sein.
\end{examplebox}

\begin{conceptbox}
Gutes maschinelles Lernen bedeutet, die richtige Balance zu finden:

\[
    \text{Modell zu einfach}
    \Rightarrow
    \text{Unteranpassung (underfitting)}.
\]

\[
    \text{Modell zu komplex}
    \Rightarrow
    \text{Überanpassung (overfitting)}.
\]
\end{conceptbox}

\section{Parametrische und nichtparametrische Modelle}

In diesem Kurs tauchen später auch parametrische und nichtparametrische Methoden auf.

\begin{definitionbox}
Ein parametrisches Modell (parametric model) hat eine feste Anzahl von Parametern.
Die Modellform wird vorher festgelegt, und das Lernen besteht darin, diese Parameter aus den Daten zu bestimmen.
\end{definitionbox}

Beispiel:

\[
    y(x,w_0,w_1) = w_0 + w_1x.
\]

Dieses Modell hat zwei Parameter, egal ob wir \(10\), \(100\) oder \(10^6\) Datenpunkte haben.

\begin{definitionbox}
Ein nichtparametrisches Modell (nonparametric model) hat nicht unbedingt eine feste Parameterzahl.
Die effektive Komplexität kann mit der Datenmenge wachsen.
\end{definitionbox}

Beispiele für nichtparametrische Methoden sind Histogramme (histograms), Kernel-Dichteschätzung (kernel density estimation) und \(k\)-nächste-Nachbarn (\(k\)-nearest neighbours).

\begin{notebox}
Nichtparametrisch bedeutet nicht, dass es gar keine Parameter gibt.
Es bedeutet eher, dass die Modellkomplexität nicht durch eine feste, kleine Parameterzahl vorgegeben ist.
\end{notebox}

\section{Warum Wahrscheinlichkeitstheorie?}

Wahrscheinlichkeitstheorie (probability theory) ist im maschinellen Lernen nicht nur ein Zusatz, sondern eine Grundlage.
Sie beantwortet Fragen wie:

\begin{itemize}
    \item Wie beschreibt man Unsicherheit mathematisch?
    \item Wie kombiniert man Vorwissen und neue Daten?
    \item Wie lernt man Parameter aus Daten?
    \item Wie entscheidet man zwischen mehreren möglichen Klassen?
    \item Wie modelliert man Rauschen (noise)?
\end{itemize}

Viele zentrale Verfahren des Kurses werden probabilistisch motiviert:

\begin{itemize}
    \item Bayes-Regel (Bayes' rule),
    \item Maximum-Likelihood-Schätzung (maximum likelihood estimation),
    \item Maximum-a-posteriori-Schätzung (maximum a-posteriori estimation),
    \item logistische Regression (logistic regression),
    \item Gaussian Mixture Models,
    \item Expectation-Maximization-Algorithmus (expectation maximization algorithm).
\end{itemize}

\begin{conceptbox}
Die probabilistische Sichtweise lautet:

\[
    \text{Lernen}
    \approx
    \text{Schließen aus unsicheren Daten}.
\]

Englisch nennt man das inference under uncertainty.
Ein Modell soll nicht nur eine Antwort geben, sondern begründen, wie plausibel verschiedene Antworten unter den Daten sind.
\end{conceptbox}

\section{Mini-Beispiel: Von Daten zu Wahrscheinlichkeiten}

Angenommen, wir haben ein einfaches Experiment mit Früchten.
Es gibt zwei Schüsseln:

\[
    s \in \{A,B\}.
\]

Aus einer Schüssel wird zufällig eine Frucht gezogen.
Die Frucht kann eine Orange oder eine weiße Beere sein:

\[
    y \in \{O,W\}.
\]

Eine typische Frage wäre:

\begin{center}
Wenn ich eine Orange beobachtet habe, wie wahrscheinlich ist es, dass sie aus Schüssel \(B\) kam?
\end{center}

Mathematisch ist das:

\[
    p(s=B \mid y=O).
\]

Das ist eine bedingte Wahrscheinlichkeit (conditional probability).
In späteren Kapiteln lernen wir, wie man solche Wahrscheinlichkeiten mit der Bayes-Regel (Bayes' rule) berechnet.

\begin{notebox}
Dieses Beispiel ist wichtig, weil es die typische ML-Situation vereinfacht zeigt:

\begin{itemize}
    \item Es gibt beobachtete Daten.
    \item Es gibt unbekannte Ursachen oder Klassen.
    \item Wir wollen aus Beobachtungen auf plausible Ursachen schließen.
\end{itemize}

Genau diese Idee steckt hinter vielen Klassifikationsmodellen.
\end{notebox}

\section{Typischer ML-Workflow}

Ein praktischer ML-Workflow sieht oft so aus:

\begin{enumerate}
    \item \textbf{Problem definieren:}
    Was soll vorhergesagt oder gelernt werden?

    \item \textbf{Daten sammeln:}
    Welche Eingaben und Zielwerte stehen zur Verfügung?

    \item \textbf{Daten vorverarbeiten:}
    Fehlende Werte, Skalierung, Kodierung, Merkmalsextraktion.

    \item \textbf{Modellklasse wählen:}
    Zum Beispiel lineare Regression, logistische Regression, neuronales Netz, PCA oder \(k\)-means.

    \item \textbf{Verlustfunktion oder probabilistisches Ziel festlegen:}
    Zum Beispiel quadratischer Fehler, negative Log-Likelihood oder ELBO.

    \item \textbf{Training durchführen:}
    Parameter aus den Trainingsdaten bestimmen.

    \item \textbf{Validieren und testen:}
    Modell auf neuen Daten bewerten.

    \item \textbf{Interpretieren:}
    Verstehen, was das Modell gut kann, wo es unsicher ist und wo es versagt.
\end{enumerate}

\begin{examtipbox}
Für Prüfungsfragen ist oft wichtig, dass du die Begriffe sauber trennst:

\begin{itemize}
    \item Modell (model): Welche Funktionsklasse oder Verteilung wird verwendet?
    \item Parameter (parameters): Welche Größen werden gelernt?
    \item Training (training): Wie werden die Parameter angepasst?
    \item Verlustfunktion (loss function): Was wird minimiert?
    \item Generalisierung (generalization): Wie gut funktioniert das Modell auf neuen Daten?
\end{itemize}
\end{examtipbox}

\section{Ausblick auf die nächsten Kapitel}

Dieses Kapitel hat die Grundideen eingeführt.
Die nächsten Kapitel machen diese Ideen mathematisch präzise.

\begin{itemize}
    \item In Kapitel 2 behandeln wir diskrete Wahrscheinlichkeiten, bedingte Wahrscheinlichkeiten, Summenregel, Produktregel und Bayes-Regel.
    \item In Kapitel 3 betrachten wir kontinuierliche Zufallsvariablen, Wahrscheinlichkeitsdichten, Erwartungswerte, Varianz, Kovarianz und Korrelation.
    \item Danach lernen wir Schätzer (estimators), Maximum-Likelihood (maximum likelihood), Maximum-a-posteriori (maximum a-posteriori) und konkrete Modelle wie lineare Regression und logistische Regression.
    \item Später folgen neuronale Netze, Backpropagation, PCA, Clustering, Gaussian Mixture Models und der EM-Algorithmus.
\end{itemize}

\section{Zusammenfassung}

\begin{itemize}
    \item Maschinelles Lernen (machine learning) bedeutet, Muster aus Daten zu lernen.
    \item Ein Modell (model) bildet Eingaben auf Ausgaben ab oder beschreibt die Struktur der Daten.
    \item Beim überwachten Lernen (supervised learning) gibt es Eingaben und Zielwerte.
    \item Beim unüberwachten Lernen (unsupervised learning) gibt es keine Zielwerte.
    \item Beim bestärkenden Lernen (reinforcement learning) lernt ein Agent durch Belohnungen.
    \item Klassifikation (classification) bedeutet Vorhersage diskreter Klassen.
    \item Regression (regression) bedeutet Vorhersage kontinuierlicher Werte.
    \item Generalisierung (generalization) ist das zentrale Ziel: Das Modell soll auf neuen Daten gut sein.
    \item Probabilistische Modelle (probabilistic models) beschreiben Unsicherheit explizit.
    \item Wahrscheinlichkeitstheorie (probability theory) ist deshalb eine zentrale Grundlage für maschinelles Lernen.
\end{itemize}

\section{Typische Prüfungsfragen}

\begin{examtipbox}
Mögliche Fragen zu diesem Kapitel:

\begin{enumerate}
    \item Erkläre den Unterschied zwischen Klassifikation (classification) und Regression (regression).
    \item Was ist der Unterschied zwischen überwachtem Lernen (supervised learning) und unüberwachtem Lernen (unsupervised learning)?
    \item Was bedeutet Generalisierung (generalization)?
    \item Warum reicht es nicht, nur den Trainingsfehler zu minimieren?
    \item Was ist Überanpassung (overfitting)?
    \item Warum sind Wahrscheinlichkeiten im maschinellen Lernen nützlich?
    \item Was ist der Unterschied zwischen Modell (model), Parameter (parameters) und Lernalgorithmus (learning algorithm)?
\end{enumerate}
\end{examtipbox}